\section{Discussion and system limitations}
While our evaluation demonstrates the soundness and efficiency of zero-knowledge spacetime compliance over finite domains, several system-level considerations govern its deployment in vehicular networks:
\begin{enumerate}
\item \textbf{Physical calibration and common-mode bias.} Cryptographic proofs guarantee computation integrity over authenticated intervals but cannot verify whether those intervals physically enclose the true event $z^*$. As illustrated by the retrospective single-receiver GNSS diagnostic in Appendix A, a static uncertainty envelope calibrated on an initial time segment (covering 76/78 epochs) covered only 2/69 epochs in a subsequent segment due to temporal drift. Practical deployments therefore require online sensor calibration, conservative quantization envelopes, and cross-provider diversity to mitigate correlated multipath or clock biases.


\item \textbf{Roster availability and gateway state management.} Requiring signatures from a fixed roster of $n$ issuers prevents provers from selectively dropping unfavorable observations, but allows a single unreachable or non-responsive issuer to cause a protocol abort. Extending the circuit with Merkle-tree roster commitments to support dynamic $n$-of-$N$ active issuer subsets, alongside distributed nonce persistence, rate limiting, and report deduplication across replicated gateways, is a natural direction for production V2X infrastructures.


\item \textbf{Privacy scope and hardware profiling.} Zero-knowledge proofs conceal the signed spacetime intervals from the gateway, but do not encrypt the report payload $R$ or prevent location inference under unrestricted adaptive policy queries. Finally, our benchmarks measure Groth16 proving latencies on an Apple M1 workstation; deploying the prover on resource-constrained automotive on-board units (OBUs) will require profiling energy consumption, V2I communication latency, and multi-party trusted setup ceremonies (or transparent SNARK backends).
\end{enumerate}


\subsection{Parameter choices and V2X interoperability}
The implemented Groth16 proving stack uses the BN254 pairing-friendly curve because the chosen gnark library supports its R1CS setup, proofs, and efficient in-circuit twisted-Edwards EdDSA/MiMC signature checks. BN254 offers approximately 100-bit security under contemporary pairing-field discrete-logarithm estimates, below a modern 128-bit long-term deployment target; these measurements therefore characterize a prototype configuration, not a 128-bit security claim. A production migration could use an appropriately selected higher-security pairing curve and remeasure setup, prover, verifier, and circuit costs. The circuit verifies gnark's twisted-Edwards BN254 EdDSA signatures, not Ed25519 or IEEE 1609.2-compliant ECDSA credentials.

In deployed C-ITS infrastructures, certificate-based V2X authenticity typically relies on IEEE 1609.2/ETSI TS 103 097 signatures, including ECDSA over NIST P-256. A practical gateway would first authenticate the V2X certificate chain and bind the authorized credential to a registered issuer identity; the in-circuit evidence signature would require either a separately enrolled SNARK-friendly key or a verified certificate-to-evidence binding. Replacing it with direct P-256 verification inside the current BN254 circuit would require non-native arithmetic and an independent cost and security evaluation. The MiMC primitive follows the pinned gnark-crypto v0.19.0 BN254/fr implementation, which uses 110 MiMC rounds, exponent five, and round constants derived from the literal seed ``seed'' through legacy Keccak-256; these parameters were not customized in this study. No compatibility, certificate-conversion, or 128-bit end-to-end performance claims are made for the evaluated prototype.

\section{Conclusion}
We have presented a zero-knowledge spacetime compliance certification framework for vehicular crowdsensing under measurement uncertainty and Byzantine evidence. Rather than certifying a single concealed coordinate, our method constructs the robust feasible set of states supported by at least $k = n - f - r$ signed uncertainty intervals and evaluates a four-valued compliance rule (Valid, Invalid, Ambiguous, and Inconsistent), accepting reports only when the feasible set is nonempty and wholly contained within the authorized policy.

We realized this relation in two exact Groth16/BN254 circuit backends---lower-corner reachability and threshold-subset enumeration---paired with a challenge-bound verification gateway. Across 12 calibrated synthetic conditions satisfying the fault and coverage budgets, the robust rule exhibited no wrong decisive classifications in the evaluated within-budget cases while improving mean decisive coverage to 87.52\% over a fixed-margin baseline (75.63\%). For small omission budgets, the subset backend achieves 164-byte proofs and 163--255 ms median proving times on an Apple M1. Future work will investigate online uncertainty calibration across heterogeneous vehicular sensors, dynamic threshold roster commitments, and end-to-end evaluation on embedded automotive hardware.